% ===========================================================================
% ΕΚΤΕΤΑΜΕΝΗ ΕΛΛΗΝΙΚΗ ΠΕΡΙΛΗΨΗ
%
% WRITE THIS LAST. It mirrors the English chapters, so every edit to them is an
% edit here. The section list below already follows the revised structure
% (verification and integration as chapters of their own), not the first
% outline - keep the two in step.
%
% SOURCE ADVANTAGE: docs/controller_design.md is already in Greek and covers
% the architecture study, the roofline and the parallelism decision. Feed it
% into §4 directly instead of translating it into English for \cref{ch:background}
% and then back again.
%
% TYPESETTING NOTE: under LGR the Greek font has no Latin glyphs, so every
% Latin run - MobileNetV2, FPGA, Vivado, int8, AXI - must be wrapped or it
% comes out as garbage. This is the single most common way this chapter breaks.
%
% The template gives you both directions and this file uses the first:
%   \selectlanguage{greek} ... \en{MobileNetV2} ...   <- whole Greek chapter
%   \textgreek{...} inside English text                <- a Greek phrase
% \en is \foreignlanguage{english}{...}, so it also gets English hyphenation.
% \textgreek was redefined \long in main.tex, so it may span paragraphs.
% ===========================================================================
\chapter*{Εκτεταμένη Ελληνική Περίληψη}
\label{ch:greek-summary}
\addcontentsline{toc}{chapter}{Εκτεταμένη Ελληνική Περίληψη}
\markboth{Εκτεταμένη Ελληνική Περίληψη}{}

\selectlanguage{greek}

\section{Εισαγωγή}
% Κίνητρο, στόχοι, συνεισφορές. Αντιστοιχεί στο \cref{ch:introduction}.
% Να αναφερθούν ρητά τα τρία επίπεδα στόχων (\en{DD-004}).

\section{Θεωρητικό Υπόβαθρο}

\subsection{Συνελικτικά Νευρωνικά Δίκτυα και \en{MobileNetV2}}
% \cref{sec:cnn}, \cref{sec:mobilenetv2}.

\subsection{Κβαντοποίηση Νευρωνικών Δικτύων}
% \cref{sec:quantization}. Το σχήμα \en{int8}, η αναδιπλωση του
% \en{batch normalization}, η επανακβαντοποίηση.

\subsection{Προγραμματιζόμενες Πύλες Πεδίου (\en{FPGAs})}
% \cref{sec:fpga}.

\subsection{Η Πλατφόρμα \en{Zynq UltraScale+ MPSoC}}
% \cref{sec:zynq}.

\section{Μοντέλο Λογισμικού Αναφοράς και Ανάλυση Φόρτου}

\subsection{Σύνολο Δεδομένων και Εκπαίδευση}
% \cref{sec:dataset}, \cref{sec:training}. Να αναφερθεί η διαρροή δεδομένων
% και ο επαναδιαχωρισμός ανά φύλλο - είναι ουσιαστικό εύρημα, όχι λεπτομέρεια.

\subsection{Κβαντοποίηση σε \en{int8}}
% \cref{sec:ptq}.

\subsection{Διανύσματα Αναφοράς και Μοντέλο σε \en{C}}
% \cref{sec:golden}, \cref{sec:c-model}. Εδώ διατυπώνεται το κεντρικό
% συμβόλαιο της εργασίας: ακρίβεια bit προς bit.

\section{Σχεδίαση του Επιταχυντή}
% Πηγή: \en{docs/controller\_design.md} - ήδη στα ελληνικά.

\subsection{Αρχιτεκτονική του Επιταχυντή}
% \cref{sec:arch-overview}.

\subsection{Παραλληλοποίηση των Μονάδων Υπολογισμού}
% \cref{sec:parallelism}. Ο πίνακας \en{32x16}, η μονάδα
% \en{depthwise}, η μονάδα \en{stem}.

\subsection{Οργάνωση Μνήμης και Μονάδα Ελέγχου}
% \cref{sec:memory}, \cref{sec:control}.

\section{Μεθοδολογία Επαλήθευσης}
% \cref{ch:verification}. ΝΕΑ ΕΝΟΤΗΤΑ σε σχέση με το αρχικό περίγραμμα, και
% πρέπει να έχει βάρος ανάλογο του αγγλικού κεφαλαίου: είναι η κύρια
% μεθοδολογική συνεισφορά.

\section{Βελτιστοποίηση Χρονισμού}
% \cref{ch:timing}. Η πορεία από τα \MHz{\thFmaxBase} στα
% \MHz{\thFmaxPhaseFour}, με το μέτρημα σε κάθε βήμα.

\section{Ολοκλήρωση με το Σύστημα Επεξεργασίας}
% \cref{ch:integration}. \en{DMA}, καταχωρητές \en{AXI-Lite},
% σχεδίαση σε επίπεδο πλακέτας, εφαρμογή στον επεξεργαστή, προεπεξεργασία.

\section{Πειραματική Αξιολόγηση}
% \cref{ch:evaluation}. Οι μετρημένοι αριθμοί και τα όριά τους - να δηλωθεί
% ότι η συχνότητα είναι ζώνη τιμών και όχι μία τιμή.

\section{Συμπεράσματα και Μελλοντικές Επεκτάσεις}
% \cref{ch:conclusion}.

\selectlanguage{english}
